% generated by scripts/make_cmame_material.py from the records; do not edit
\begin{table}[tbp]
\centering
\footnotesize
\setlength{\tabcolsep}{4pt}
\caption{Post hoc, three-dimensional: the energy-optimal amplitude $c_b$ of the four load cases of an instance together, on fixed geometries and loads meshed at four mesh sizes $l_c$.}
\label{tab:remesh}
\resizebox{\ifdim\width>\textwidth\textwidth\else\width\fi}{!}{%
\begin{tabular}{lcccc}
\toprule
Network & $l_c$ = 0.0906 & $l_c$ = 0.0742 & $l_c$ = 0.0579 & \shortstack{$l_c$ = 0.0374 \\ (fine)} \\
\midrule
Transformer & 0.98 $\pm$ 0.00 & 1.00 $\pm$ 0.01 & 1.01 $\pm$ 0.00 & 1.27 $\pm$ 0.08 \\
Bottleneck, M = 512 & 0.91 $\pm$ 0.04 & 1.00 $\pm$ 0.02 & 1.14 $\pm$ 0.09 & 1.96 $\pm$ 0.46 \\
Bottleneck, M = 1,024 & 0.94 $\pm$ 0.01 & 1.01 $\pm$ 0.01 & 1.14 $\pm$ 0.02 & 1.99 $\pm$ 0.35 \\
\bottomrule
\end{tabular}}

\smallskip
\begin{minipage}{\linewidth}\footnotesize\raggedright
16 fresh geometries, each meshed at the 4 $l_c$ with identical geometry and loads, so only the mesh changes; every mesh evaluated by every seed. $c_b$: the single energy-optimal amplitude of all four load cases of an instance (label-free, \cref{sec:transfer:amplitude}); median over the geometries, then seed mean $\pm$ sample SD. Training range: $l_c$ = 0.0579--0.0906.\par
The largest nodal load over the four load cases, by which both networks multiply their output, relative to $l_c$ = 0.0906, at $l_c$ = 0.0742 / 0.0579 / 0.0374: 0.705 / 0.430 / 0.186; for comparison ($l_c$ / 0.0906)$^2$: 0.671 / 0.408 / 0.170.\par
Predicted energy norm relative to $l_c$ = 0.0906 (label-free), seed means, $l_c$ = 0.0742 / 0.0579 / 0.0374: transformer: 0.99 / 0.98 / 0.75; M = 512: 0.90 / 0.82 / 0.50; M = 1,024: 0.93 / 0.83 / 0.49.\par
Measured after the runs.
\end{minipage}
\end{table}
